% Generated draft: requires mathematical and source review.
\documentclass[11pt]{article}
\usepackage[margin=25mm]{geometry}
\usepackage{fontspec}
\setmainfont{Latin Modern Roman}
\usepackage{amsmath,amssymb,mathtools,unicode-math}
\setmathfont{Latin Modern Math}
\usepackage{enumitem,needspace,xurl,hyperref}
\setlength{\parindent}{0pt}
\setlength{\parskip}{5pt}
\setlength{\emergencystretch}{3em}
\newcommand{\N}{\mathbb{N}}
\newcommand{\Z}{\mathbb{Z}}
\newcommand{\Q}{\mathbb{Q}}
\newcommand{\R}{\mathbb{R}}
\newcommand{\C}{\mathbb{C}}
\newcommand{\F}{\mathbb{F}}
\newcommand{\A}{\mathbb{A}}
\newcommand{\PP}{\mathbb{P}}
\newcommand{\E}{\mathbb{E}}
\newcommand{\eps}{\varepsilon}
\newcommand{\dd}{\,\mathrm{d}}
\begin{document}
\section*{UnsolvedMath ID: 30004496}
\textbf{Source problem code:} catalogue entry
\subsection{Definitions and mathematical statement}
Let \(M\) be a compact arithmetic hyperbolic \(3\)-manifold, and let \(d\mu\) denote the volume measure intended in the source, whose normalization is unspecified. A Hecke--Laplace eigenfunction is a nonzero smooth function \(f:M\to\C\) that is an eigenfunction of the Laplace operator and a simultaneous eigenfunction of the relevant Hecke operators. Writing the Laplace operator with the nonnegative sign convention, this means
\[
\Delta_M f=\lambda_f f,\qquad T f=a_T f\quad(T\in\mathcal H),
\]
where \(\mathcal H\) denotes the intended family of Hecke operators. The arithmetic data and the family \(\mathcal H\) are not specified in the supplied record. Assume the stated normalization
\[
\|f\|_{L^2(M,\mu)}^2=\int_M |f(x)|^2\,d\mu(x)=1.
\]
Define the maximum magnitude and root-mean-square magnitude by
\[
\|f\|_\infty=\max_{x\in M}|f(x)|,
\qquad
\operatorname{rms}_\mu(f)=\left(\frac{1}{\mu(M)}\int_M|f(x)|^2\,d\mu(x)\right)^{1/2}
=\mu(M)^{-1/2}.
\]
Thus the root-mean-square magnitude equals \(1\) if \(\mu(M)=1\). The stated integral normalization alone does not identify the average magnitude \(\mu(M)^{-1}\int_M|f|\,d\mu\) with \(1\).

\textbf{Source question, with unresolved qualifications.} Must \(|f|\) attain values much larger than its root-mean-square magnitude? Equivalently, the qualitative question concerns the size of the dimensionless ratio
\[
\frac{\|f\|_\infty}{\operatorname{rms}_\mu(f)}
=\mu(M)^{1/2}\|f\|_\infty.
\]
The supplied wording does not define ``much larger'': it specifies neither a numerical threshold nor an asymptotic lower bound, a limiting parameter, or quantifiers over eigenfunctions. It also does not specify whether \(M\) is fixed, whether a positive-eigenvalue restriction is intended, or whether the assertion should concern every eigenfunction, a subsequence, or some other family. Consequently, the record supplies a qualitative question rather than a uniquely determined quantitative conjecture. No particular growth assertion is imposed here.
\subsection{Short English statement}
The source asks whether an \(L^2\)-normalized Hecke--Laplace eigenfunction on a compact arithmetic hyperbolic three-manifold must have unusually large peak values. The comparison size is \(1\) only when the volume measure has total mass \(1\). The meaning of ``much larger'' and the intended family and quantifiers remain unspecified.
\subsection{Sources}
\begin{itemize}
\item Local UnsolvedMath registry: \texttt{data/unsolvedmath/problems.json}, permanent ID \(30004496\), title ``Large Values of Arithmetic Hyperbolic Three-Manifold Eigenfunctions,'' source problem code \texttt{catalogue entry}. This draft uses the supplied complete source record.
\item Supplied report citation: \emph{Automorphic Forms and Arithmetic (2021)}, Oberwolfach Reports, DOI \texttt{10.4171/owr/2020/26}, pp.~1317--1383. Supplied source and literature link, labeled ``Oberwolfach Report 26/2020'': \url{https://doi.org/10.4171/owr/2020/26}. The record attributes the question to the Oberwolfach Reports workshop collection; no individual proposer is supplied. The proposed year is recorded as \(2020\). The citation's \(2021\) title date and its \(2020\) report designation are retained as supplied.
\item Literature assessment dated \texttt{2026-08-22}: the supplied curation labels the question open, describes it as a qualitative large-values question, and states that the extracted wording supplies no quantitative threshold that has been proved. This is supplied dated curation, not an independently verified literature assessment; the record requests independent verification of time-sensitive status and cited claims.
\item The registry labels the statement \texttt{corrected\_verified} and describes its editing as preserving mathematical content while repairing presentation or omitted definitions. This is a supplied editorial label, not independent mathematical verification. The original wording compares with an ``average value 1,'' whereas the cleaned wording compares with a ``root-mean-square size 1.'' Their measure-normalization qualification remains unresolved.
\end{itemize}
\end{document}
